\documentclass[11pt,a4paper]{article}
\usepackage[margin=1in]{geometry}
\usepackage{amsmath,amssymb,amsfonts}
\usepackage{booktabs}
\usepackage{hyperref}
\usepackage{xcolor}
\usepackage{tcolorbox}

\newcommand{\vE}{\mathbf{E}}
\newcommand{\vB}{\mathbf{B}}

\title{Which $\kappa$ Is Correct?\\
Resolving the Constitutive Split Parameter\\
in Density Field Dynamics}

\author{Analysis for G.\ T.\ Alcock --- DFD Research Programme\\
March 23, 2026}
\date{}

\begin{document}
\maketitle

%=============================================================================
\section{The Problem}
%=============================================================================

The vacuum loading paper introduces a constitutive split parameter $\kappa$
controlling how the electric and magnetic sectors respond asymmetrically
to the $\psi$-field:
\begin{equation}\label{eq:constitutive}
\varepsilon(\psi) = \varepsilon_0\, n\, e^{+\kappa\psi}, \qquad
\mu_{\rm mag}(\psi) = \mu_0\, n\, e^{-\kappa\psi},
\end{equation}
with $n = e^\psi$.  The product $\varepsilon\mu = \varepsilon_0\mu_0\,n^2$
is independent of $\kappa$, so the phase velocity $v_{\rm ph} = c/n$ is
preserved regardless of $\kappa$.

Two numerical values appear in the Deep Physics extension paper:
\begin{center}
\begin{tabular}{lcc}
\toprule
\textbf{Route} & \textbf{Value} & \textbf{Source} \\
\midrule
Euler--Heisenberg (EH) & $\kappa_{\rm EH} = 3/11 \approx 0.273$ &
Ratio of EH coefficients $4:7$ \\
Electroweak mixing & $\kappa_{\rm EW} = \cos 2\theta_W \approx 0.538$ &
$W^3/B$ mixing angle \\
\bottomrule
\end{tabular}
\end{center}

These differ by a factor of $\sim 2$.  The paper favors the electroweak
value but does not definitively resolve the discrepancy.  This document
determines which is correct, whether they apply at different scales, or
whether neither is right.

%=============================================================================
\section{What the DFD Action Actually Says}
\label{sec:action}
%=============================================================================

\subsection{The minimal DFD action predicts $\kappa = 0$}

The core DFD Lagrangian for the EM--$\psi$ sector is:
\begin{equation}\label{eq:L-DFD}
\mathcal{L}_{\rm EM} = -\frac{1}{4}e^{2\psi}\,F_{\mu\nu}F^{\mu\nu}.
\end{equation}
The factor $e^{2\psi}$ multiplies the \emph{full} $F^2 = F_{\mu\nu}F^{\mu\nu}
= 2(\vB^2 - \vE^2/c^2)$ without any $\vE$/$\vB$ asymmetry.  Through the
Tamm--Plebanski construction, this optical metric yields:
\begin{equation}
\varepsilon_{\rm eff} = \varepsilon_0\, e^{+\psi}, \qquad
\mu_{\rm eff} = \mu_0\, e^{+\psi}.
\end{equation}
This is the \textbf{symmetric case}: $\kappa = 0$.  The electric and
magnetic sectors respond identically.  This is confirmed explicitly in
Section~10 of the main DFD paper (cavity--atom analysis), where the
constitutive relations are written with $\varepsilon_{\rm eff}$ and
$\mu_{\rm eff}$ both scaling as $e^{+\psi}$, with no split.

\begin{tcolorbox}[colback=yellow!5!white, colframe=yellow!60!black,
title=Key Point 1: The Minimal DFD Action Gives $\kappa = 0$]
The classical DFD Lagrangian $\mathcal{L} \supset -\frac{1}{4}e^{2\psi}F^2$
treats $\vE$ and $\vB$ symmetrically.  Any $\kappa \neq 0$ must arise from
\textbf{quantum corrections}---it is not a tree-level prediction of the
DFD postulates.
\end{tcolorbox}

\subsection{Why this matters}

The tree-level $\kappa = 0$ result means that $\kappa$ is a \emph{radiative
correction} to the DFD vacuum.  The question ``which $\kappa$ is correct?''
therefore becomes: ``which quantum effect generates the leading $\varepsilon/\mu$
asymmetry in the DFD vacuum?''  This reframing is essential because it tells
us that $\kappa$ must be proportional to loop factors ($\alpha/\pi$ or
similar), which constrains its magnitude.

%=============================================================================
\section{Analysis of the Two Derivation Routes}
\label{sec:analysis}
%=============================================================================

\subsection{Route 1: Euler--Heisenberg ($\kappa_{\rm EH} = 3/11$)}

\subsubsection{The derivation}

The EH effective Lagrangian at one loop is:
\begin{equation}
\mathcal{L}_{\rm EH} = \frac{\alpha}{90\pi}\frac{1}{B_c^2}
\left[4(\vE^2 - \vB^2)^2 + 7(\vE\cdot\vB)^2\right].
\end{equation}
The Deep Physics paper extracts a ``split'' from the coefficient ratio:
$\kappa_{\rm EH} = (7-4)/(7+4) = 3/11$.

\subsubsection{What this actually measures}

The EH Lagrangian describes \textbf{photon--photon scattering} mediated by
virtual electron loops.  The coefficients $4$ and $7$ multiply the two
independent Lorentz-invariant combinations of $F_{\mu\nu}$:
\begin{itemize}
\item The $\mathcal{F}^2 = (\vE^2 - \vB^2)^2$ term: photon self-interaction
  in the scalar channel.
\item The $\mathcal{G}^2 = (\vE\cdot\vB)^2$ term: photon self-interaction
  in the pseudoscalar channel.
\end{itemize}

\textbf{Neither of these is the $\varepsilon/\mu$ split under gravitational
loading.}  The EH Lagrangian tells us how photons scatter off each other,
not how the vacuum's permittivity and permeability respond differently to a
background $\psi$-field.  The ``3/11'' is the asymmetry between two
\emph{photon--photon} scattering channels, not between the electric and
magnetic vacuum polarization responses to $\psi$.

\subsubsection{Validity regime}

The EH effective Lagrangian is valid for photon energies
$\omega \ll m_e c^2$ and field strengths $E, B \ll B_c = m_e^2 c^3/(e\hbar)
\approx 4.4 \times 10^9$~T.  These conditions are satisfied in all
laboratory and astrophysical contexts relevant to DFD.  However, the
validity of the EH Lagrangian is not the issue---the issue is that it
computes the wrong quantity.

\begin{tcolorbox}[colback=red!5!white, colframe=red!60!black,
title=Verdict on EH Route]
\textbf{The EH derivation of $\kappa = 3/11$ is incorrect.}  It conflates
the photon--photon scattering asymmetry ($\mathcal{F}^2$ vs.\
$\mathcal{G}^2$ channels) with the constitutive $\varepsilon/\mu$ split
under $\psi$-loading.  These are different physical quantities.  The
$4:7$ ratio has no direct bearing on how $\varepsilon$ and $\mu$
respond to a slowly varying background scalar field.
\end{tcolorbox}

\subsection{Route 2: Electroweak mixing ($\kappa_{\rm EW} = \cos 2\theta_W$)}

\subsubsection{The derivation}

The physical photon is a mixture of the hypercharge gauge boson $B$ and
the third isospin component $W^3$:
\begin{equation}
A_\mu = B_\mu \cos\theta_W + W^3_\mu \sin\theta_W.
\end{equation}
The Deep Physics paper argues that the $\psi$-coupling distinguishes
between these two components, giving:
\begin{equation}
\kappa = \frac{\cos^2\theta_W - \sin^2\theta_W}
{\cos^2\theta_W + \sin^2\theta_W} = \cos 2\theta_W \approx 0.538.
\end{equation}

\subsubsection{What this actually measures}

This derivation assumes that the $B$ and $W^3$ components of the photon
couple to $\psi$ with different strengths proportional to $\cos^2\theta_W$
and $\sin^2\theta_W$ respectively, and that this difference maps onto an
$\varepsilon/\mu$ asymmetry.

There are two problems:
\begin{enumerate}
\item \textbf{Why would $B$ vs.\ $W^3$ map onto $\varepsilon$ vs.\ $\mu$?}
  The decomposition of the photon into $B$ and $W^3$ is an internal gauge
  structure, not a spatial $\vE/\vB$ decomposition.  The hypercharge and
  isospin components both contribute to $\vE$ and $\vB$ equally---they do
  not preferentially couple to one Lorentz sector over the other.

\item \textbf{The coupling to $\psi$ is gravitational, not gauge-selective.}
  In DFD, $\psi$ couples through the stress-energy tensor
  $T_{\mu\nu}^{\rm EM}$, which is the \emph{total} EM stress-energy.
  Both $B$ and $W^3$ contribute to the same $T_{\mu\nu}$ after symmetry
  breaking.  The $\psi$-field does not see the internal gauge decomposition;
  it sees the physical EM fields.
\end{enumerate}

\subsubsection{When the electroweak argument could work}

The argument would be valid if $\psi$ coupled not through $T_{\mu\nu}$ but
through a gauge-specific channel---for instance, if $\psi$ were sourced
by the \emph{hypercharge} current rather than the gravitational stress
tensor.  In such a model, the $B$-component of the photon would couple
more strongly than the $W^3$ component, and the mixing angle would enter.
But this contradicts the DFD postulates, where $\psi$ is the gravitational
potential and couples universally through stress-energy.

\begin{tcolorbox}[colback=orange!5!white, colframe=orange!60!black,
title=Verdict on Electroweak Route]
\textbf{The electroweak derivation of $\kappa = \cos 2\theta_W$ has a
stronger physical motivation than the EH route, but it relies on an
assumption that is not justified within the DFD framework}: that the
internal $B/W^3$ decomposition of the photon maps onto a spatial
$\varepsilon/\mu$ asymmetry.  Since $\psi$ couples gravitationally
(through $T_{\mu\nu}$), it does not distinguish between the gauge
components of the physical photon.

However, this route \textbf{could become correct} if the DFD internal
space ($\mathbb{C}P^2 \times S^3$) generates a gauge-selective $\psi$
coupling at one loop.  This would require an explicit spectral-action
computation.
\end{tcolorbox}

%=============================================================================
\section{The Correct Framework: Vacuum Polarization Tensor Decomposition}
\label{sec:correct}
%=============================================================================

\subsection{What actually generates $\kappa \neq 0$}

The constitutive split $\kappa$ measures how the vacuum's \emph{dielectric
response} differs between the electric and magnetic channels when the vacuum
is loaded by $\psi$.  The correct computation requires:

\begin{enumerate}
\item Start from the DFD action $\mathcal{L} = -\frac{1}{4}e^{2\psi}F^2$
  in a background $\psi$-field.
\item Compute the one-loop vacuum polarization tensor
  $\Pi^{\mu\nu}(q; \psi)$ for virtual charged fermion loops in this
  background.
\item Decompose $\Pi^{\mu\nu}$ into electric ($E$-type) and magnetic
  ($B$-type) projections in the rest frame of the $\psi$-gradient:
  \begin{equation}
  \Pi^{\mu\nu} = \Pi_E(q^2;\psi)\,P_E^{\mu\nu}
  + \Pi_B(q^2;\psi)\,P_B^{\mu\nu}.
  \end{equation}
\item The constitutive split is then:
  \begin{equation}\label{eq:kappa-correct}
  \kappa = \frac{1}{2}\frac{\partial}{\partial\psi}
  \ln\frac{\Pi_E(0;\psi)}{\Pi_B(0;\psi)}\bigg|_{\psi=0}.
  \end{equation}
\end{enumerate}

\subsection{What the computation gives}

In standard QED (which is the relevant one-loop computation), the vacuum
polarization in an external gravitational field has been computed.  The
key result is from the Drummond--Hathrell effect (1980): photon
propagation in curved spacetime acquires corrections of the form
\begin{equation}
\Delta\mathcal{L} = \frac{\alpha}{90\pi m_e^2}
\left(5 R_{\mu\nu\alpha\beta}F^{\mu\nu}F^{\alpha\beta}
- 26 R_{\mu\nu}F^{\mu\alpha}F^\nu{}_\alpha
+ 2R\,F_{\mu\nu}F^{\mu\nu}\right).
\end{equation}

For a static, spherically symmetric background (the DFD case, where the
Riemann tensor is determined by $\psi$), this produces a polarization-dependent
correction to the photon dispersion relation.  The $\vE$ and $\vB$ modes
propagate at slightly different speeds---this is gravitational birefringence.

The magnitude of this effect is:
\begin{equation}
\frac{\delta v_E - \delta v_B}{c} \sim \frac{\alpha}{m_e^2}\,R
\sim \frac{\alpha}{m_e^2}\,\frac{GM}{r^3},
\end{equation}
which is suppressed by $\alpha$ \emph{and} by the ratio of the Compton
wavelength to the curvature radius.  For any astrophysical or laboratory
setting, this is negligibly small ($\lesssim 10^{-40}$).

\subsection{Implication for $\kappa$}

The Drummond--Hathrell effect tells us that $\kappa$ from QED vacuum
polarization is not a number of order unity.  It is:
\begin{equation}\label{eq:kappa-QED}
\kappa_{\rm QED} \sim \frac{\alpha}{90\pi}\,\frac{\hbar^2}{m_e^2 c^2 L^2}
\sim 10^{-44}\left(\frac{1~\text{m}}{L}\right)^2,
\end{equation}
where $L$ is the curvature scale of the $\psi$-background.  For any
macroscopic curvature scale, $\kappa_{\rm QED}$ is absurdly small.

\begin{tcolorbox}[colback=green!5!white, colframe=green!60!black,
title=Key Point 2: QED Predicts $\kappa \approx 0$]
The one-loop QED vacuum polarization in a curved/loaded background
generates $\kappa \sim \alpha \times (\lambdabar_e/L)^2$, which is
unmeasurably small for macroscopic $\psi$-gradients.  Both the EH value
($3/11$) and the electroweak value ($\cos 2\theta_W$) are order unity,
which is \textbf{incompatible} with the actual loop computation.

The correct QED result is that $\kappa$ is effectively zero at
laboratory and astrophysical scales.
\end{tcolorbox}

%=============================================================================
\section{Resolution: Three Distinct Quantities}
\label{sec:resolution}
%=============================================================================

The confusion arises because three physically distinct quantities have
been conflated under the single symbol $\kappa$:

\begin{center}
\renewcommand{\arraystretch}{1.4}
\begin{tabular}{p{3cm}p{3cm}p{3cm}p{3cm}}
\toprule
\textbf{Quantity} & \textbf{Value} & \textbf{What it measures} &
\textbf{Relevance to DFD} \\
\midrule
EH photon--photon asymmetry &
$\frac{7-4}{7+4} = \frac{3}{11}$ &
Ratio of $\mathcal{G}^2$ to $\mathcal{F}^2$ scattering amplitudes
in the EH effective action &
None.  This describes photon self-interaction, not the vacuum's
response to $\psi$.  \\
\midrule
Electroweak $B/W^3$ mixing &
$\cos 2\theta_W \approx 0.54$ &
Asymmetry between hypercharge and isospin components of the photon &
Indirect at best.  Maps onto $\varepsilon/\mu$ only if $\psi$ couples
gauge-selectively, which contradicts universal gravitational coupling. \\
\midrule
QED vacuum birefringence in curved background &
$\sim\frac{\alpha}{90\pi}\frac{\lambdabar_e^2}{L^2} \approx 0$ &
Actual $E/B$ polarization asymmetry from vacuum polarization in a
$\psi$-background &
Direct.  This is the correct computation.  The result is effectively
zero. \\
\bottomrule
\end{tabular}
\end{center}

%=============================================================================
\section{Does $\kappa$ Run with Energy?}
\label{sec:running}
%=============================================================================

Since the two proposed values ($3/11$ and $\cos 2\theta_W$) differ by a
factor of $\sim 2$, one might ask whether $\kappa$ runs with energy scale,
with one value applying in the infrared and the other at $M_Z$.

\subsection{Running of $\sin^2\theta_W$}

The weak mixing angle runs from $\sin^2\theta_W(M_Z) = 0.2312$ to
$\sin^2\theta_W(0) \approx 0.238$ (MS-bar scheme).  This gives:
\begin{equation}
\cos 2\theta_W(0) = 1 - 2(0.238) = 0.524, \qquad
\cos 2\theta_W(M_Z) = 0.538.
\end{equation}
The variation is only $\sim 3\%$ across the full energy range.  There is
no energy at which $\cos 2\theta_W$ approaches $3/11 = 0.273$.

\subsection{Can the EH and EW values agree at some scale?}

For $\cos 2\theta_W(\mu) = 3/11$, we would need $\sin^2\theta_W(\mu)
= (1 - 3/11)/2 = 4/11 \approx 0.364$.  This does not occur at any
perturbatively accessible scale.  It would require $\sin^2\theta_W$ to
increase by more than $50\%$ from its $M_Z$ value, which would require
running well beyond the perturbative regime into new physics.

\begin{tcolorbox}[colback=blue!5!white, colframe=blue!50!black,
title=Verdict on Running]
The two values do \textbf{not} correspond to the same quantity at different
scales.  $\cos 2\theta_W$ varies by only $\sim 3\%$ across the entire
accessible energy range, and never approaches $3/11$.  The discrepancy
is not resolved by RG running.
\end{tcolorbox}

%=============================================================================
\section{Recommendations for the DFD Programme}
\label{sec:recommendations}
%=============================================================================

\subsection{What should the vacuum loading paper say?}

\begin{enumerate}
\item \textbf{The minimal DFD action predicts $\kappa = 0$ at tree level.}
  This should be stated clearly.  The constitutive relations from the
  Tamm--Plebanski construction with $n = e^\psi$ are symmetric:
  $\varepsilon_{\rm eff} = \varepsilon_0 e^{\psi}$,
  $\mu_{\rm eff} = \mu_0 e^{\psi}$.

\item \textbf{$\kappa \neq 0$ is a radiative correction.}  It arises from
  the one-loop QED (or electroweak) vacuum polarization in the $\psi$
  background.  Its magnitude is suppressed by loop factors and by the
  ratio of the Compton wavelength to the curvature scale.

\item \textbf{The EH derivation should be withdrawn.}  The $3/11$ ratio
  measures photon--photon scattering channel asymmetry, not the
  $\varepsilon/\mu$ split under gravitational loading.

\item \textbf{The electroweak derivation needs reformulation.}  The
  $\cos 2\theta_W$ value has a suggestive structure but requires an
  explicit mechanism by which the $B/W^3$ decomposition maps onto the
  spatial $\vE/\vB$ split.  This mechanism does not exist within the
  standard gravitational ($T_{\mu\nu}$) coupling.  It \emph{could}
  exist if the $\mathbb{C}P^2 \times S^3$ spectral action generates a
  gauge-selective $\psi$ coupling at one loop.  This would be a genuine
  prediction but requires explicit computation.

\item \textbf{Present $\kappa$ as a phenomenological parameter.}  Until
  the spectral-action computation is done, $\kappa$ should be treated as
  a free parameter constrained by experiment ($|\kappa| \lesssim 1$ from
  cavity stability).  The two ``predicted'' values ($3/11$ and $\cos 2\theta_W$)
  should be presented as \emph{motivations} for experimental searches,
  not as DFD predictions.
\end{enumerate}

\subsection{What computation would settle the question?}

The definitive answer requires computing the $a_4$ Seeley--DeWitt coefficient
for the EM sector on the Toeplitz-truncated $\mathbb{C}P^2 \times S^3$
geometry, decomposed into $E$-type and $B$-type projections.  This is the
same spectral-action machinery used for the $\alpha$ derivation, but
applied to the \emph{split} between electric and magnetic stiffnesses
rather than their sum.

Specifically:
\begin{equation}
\kappa = \frac{\kappa_{E} - \kappa_{B}}{\kappa_{E} + \kappa_{B}},
\end{equation}
where $\kappa_{E}$ and $\kappa_{B}$ are the gauge coupling stiffnesses
for the $E$-type and $B$-type polarization channels computed from the
spectral action.  If the $\mathbb{C}P^2 \times S^3$ geometry treats
these channels identically (which it does at leading order, since the
geometry is isotropic), then $\kappa = 0$ exactly.  If the Toeplitz
truncation or the hypercharge structure breaks this symmetry, then
$\kappa$ is calculable and nonzero.

\subsection{Experimental implications}

Regardless of which theoretical value is correct, the experimental
program is unchanged:
\begin{itemize}
\item TE/TM cavity comparisons measure $\kappa\psi$ directly.
\item Multi-species clock comparisons constrain $\kappa$ through
  species-dependent LPI residuals.
\item A null result ($\kappa = 0$ to experimental precision) would confirm
  the minimal DFD action and rule out both proposed values.
\item A detection of $\kappa \neq 0$ would identify new physics beyond
  the tree-level DFD framework.
\end{itemize}

%=============================================================================
\section{Summary}
\label{sec:summary}
%=============================================================================

\begin{tcolorbox}[colback=blue!5!white, colframe=blue!50!black,
title=Final Verdict]
\textbf{Neither value is correct as a tree-level DFD prediction.}

\begin{enumerate}
\item The \textbf{minimal DFD action} ($-\frac{1}{4}e^{2\psi}F^2$) predicts
$\kappa = 0$ at tree level.  Both $\varepsilon$ and $\mu$ scale as
$e^{+\psi}$ symmetrically.

\item The \textbf{Euler--Heisenberg value} ($3/11$) is wrong: it measures
photon--photon scattering asymmetry, not the $\varepsilon/\mu$ response
to $\psi$-loading.

\item The \textbf{electroweak value} ($\cos 2\theta_W \approx 0.54$) has
the right flavor but lacks a rigorous mapping from the internal $B/W^3$
decomposition to the spatial $\varepsilon/\mu$ split.

\item The \textbf{QED vacuum polarization} in a curved background gives
$\kappa \sim \alpha \times (\lambdabar_e/L)^2 \approx 0$, consistent
with $\kappa = 0$ at all accessible scales.

\item The two values do \textbf{not} correspond to different energy scales;
$\cos 2\theta_W$ never runs to $3/11$ at any perturbative scale.

\item The definitive answer requires the \textbf{spectral-action
computation} on $\mathbb{C}P^2 \times S^3$ with $E/B$ polarization
decomposition.  Until this is done, $\kappa$ should be treated as a
\textbf{free phenomenological parameter} constrained by experiment.
\end{enumerate}
\end{tcolorbox}

\end{document}
